\section{How to use this document}
\label{sec:howto}

This document has two audiences and it is worth saying which paragraph is
for whom.

A reader who wants the mathematics should read Part~\ref{part:exact},
which is self-contained and compiles on its own as a separate paper
(\srcpath{paper/master/workhouse_band_paper.tex}). Nothing in
Part~\ref{part:exact} depends on anything later, and nothing in it is
conditional on an open estimate.

A reader who intends to \emph{work} on the continuum problem should read
\S\ref{sec:breaks}, \S\ref{sec:register} and \S\ref{sec:obligations}
first, in that order, and only then come back for the mathematics. The
reason is economic. This program has spent four years discovering that
most of the obvious routes to the continuum limit are closed, and the
record of \emph{which} routes are closed and \emph{why} is worth more
than the record of what is open. Four sessions of one earlier effort were
spent on a single route that a static argument had already shown could
not decide the question it was aimed at; that experience is why the
route register in \S\ref{sec:obligations} exists and why dead routes are
reproduced in it at full length rather than deleted.

\subsection{What the badges mean}

Every statement in this document carries a machine-verification badge.

\begin{itemize}
  \item \Tzero{} --- Lean~4 compiles the statement from explicit
    definitions with no \texttt{sorry}, using only \texttt{propext},
    \texttt{Classical.choice} and \texttt{Quot.sound}. Nothing written
    in prose can weaken a T0 badge.
  \item \Tone{} --- the statement is re-derived symbolically, in exact
    arithmetic, from stated definitions.
  \item \Ttwo{} --- floating-point agreement within a tolerance that the
    check itself prints.
  \item \Tthree{} --- no single repository check covers the whole
    statement.
\end{itemize}

The badge is \emph{independent} of mathematical status. A theorem with a
complete analytic proof and no formalization carries \Tthree{}: the badge
says what the machine checks, not what the mathematics establishes. This
is the one convention a reader must internalize to read the document
correctly, because the opposite reading --- T3 as a synonym for
``unproved'' --- makes most of Part~\ref{part:continuum} look weaker than
it is, and the reading of T0 as a synonym for ``important'' makes it look
stronger. Both errors have been made in this program's own internal
documents and corrected.

Where a statement's scope is narrower than its headline suggests, a
\textbf{Does not assert} paragraph follows it. Those paragraphs are not
hedging. They are the part of each result that took longest to get right,
and several of them exist because an earlier version of the claim was
broader and was refuted.

\subsection{If you are the next agent on this program}

The following is an operating procedure, not advice. It is written for
the case --- now the normal case --- in which the next piece of work on
this program is done by an automated agent with a fresh context window.

\begin{enumerate}
  \item \textbf{Do not trust this document over the ledgers.} Every
    number quoted in the prose here is generated from
    \srcpath{ledger/*.yaml} by \srcpath{paper/master/generate\_apparatus.py}.
    If the ledger and this document disagree, the ledger is right and
    this document is stale; regenerate it. The counts in
    \S\ref{sec:obligations}, \S\ref{sec:external} and
    Appendix~\ref{app:verification} are macros, not typed literals,
    precisely so that this failure mode is visible rather than silent.

  \item \textbf{Establish the state before choosing a target.} Run
    \texttt{workhouse brief -{}-startup}, then \texttt{workhouse brief <gap> -{}-json}
    for the gap you mean to work on, and retain the snapshot. A saved
    graph in a fresh clone reports freshness \texttt{unknown} until local
    generation provenance exists; saved briefing mode never silently
    rebuilds it.

  \item \textbf{Read the dead routes for your gap before starting.}
    \texttt{workhouse why <gap>} prints them. \S\ref{sec:obligations}
    reproduces them here. Of \RouteTotal{} recorded routes across all
    gaps, \RouteDead{} are dead, and several died only after substantial
    expenditure with a measured cost recorded in the ledger. A route
    marked \texttt{cannot\_decide} is not merely unpromising: it is
    structurally incapable of resolving the claim it was aimed at, and
    re-running it will produce output that answers a different question.

  \item \textbf{Do not aim at the continuum limit directly.}
    \S\ref{sec:breaks} states six independent obstructions, five of them
    with explicit finite counterexamples. Each is a constraint on what a
    correct argument can look like, and an argument that does not
    explicitly clear all six is not a candidate. In particular, Break~5
    (\S\ref{sec:break5}) shows that the regime in which this program's
    strongest construction holds and the regime in which the continuum
    limit lives are disjoint sets of couplings. That is a statement about
    the shape of any possible proof, not a gap in the present one.

  \item \textbf{Pick a numbered Problem from \S\ref{sec:obligations}}, not
    a topic. Each Problem states its available hypotheses, what it
    unblocks downstream, and the specific next computation. Problems are
    ordered by the ledger's research judgement, which is not a claim that
    earlier ones are prerequisites for later ones.

  \item \textbf{Record the attempt, including its failure.} A route that
    dies is a result. Set its state in \srcpath{ledger/gaps.yaml}, name
    what closed it in \texttt{closed\_by}, and if the instrument could
    not have reached the claim, say so in \texttt{cannot\_decide}. The
    working agreement in \srcpath{CLAUDE.md} is the binding version of
    this instruction; this paragraph only explains why it is there.
\end{enumerate}

\subsection{What would falsify the contents of this document}

Three things, in decreasing order of how much they would cost.

First, an error in the Haar or Weingarten input to \S\ref{sec:second}
would propagate to every order and to both the fourth-order band and its
rank grading. That input is now representation theory rather than
assumption (\S\ref{sec:second}), and it is the load-bearing algebraic
fact of Part~\ref{part:exact}.

Second, a demonstration that the plaquette source family of
\S\ref{sec:wilson} is not total in the physical complement would empty
the fixed-spacing band theorem of physical content without making it
false --- this is exactly the failure mode that Break~4
(\S\ref{sec:break4}) exhibits in a three-dimensional model, and totality
is therefore asserted in that theorem as a proved conclusion rather than
assumed.

Third, and least costly, a published result reproducing any of the exact
coefficients of \S\ref{sec:grading} or \S\ref{sec:fourth} would remove
this program's claim to have derived them first without touching their
correctness. \S\ref{sec:external} records the searches already performed
against \LitPapers{} indexed papers.
